\cotefile{134-2}
\scanpage{154}\mspage{9}
3) On peut faire les ${=}\mathrm{s}$ choses pour des coeff tels que $\mathbb{Z}_{\unsure{\ell|e}}$ (Jouanolou, thèse -- non publiée je crains !)

\inserted{4) Ça marche pour $f:X\longrightarrow S$ prés. finie séparé, si $f$ a des propriétés de ``locale trivialité cohomologique'' (propr. locales \uline{en haut}) p.ex. $f$ \uline{lisse} ; on trouve $\underline{H}_{*X/S}=Rf_{!}\,f^{!}(\Lambda_S)$.}

B) Artin et Mazur ont étudié, dans un esprit voisin du tien, \uline{l'autodualité} de la jacobienne d'une courbe relative $X/S$. Il faut les consulter pour les résultats précis, peut-être cela marche dès que $X/S$ est propre et plat de dim relative $1$ -- en tous cas c'est OK sur un anneau de valuation discrète à fibre \uline{générique} lisse. La fibre spéciale peut être très sauvage. (J'ai utilisé \deleted{\unsure{pour|leur}}\unsure{leur|leurs} résultat dans SGA \deleted{\unsure{3}}7 pour prouver, dans le cas des jacobiennes, une conjecture de dualité sur les groupes de composantes connexes associés aux \deleted{\unsure{de|des} V.A.} modèles de Néron de V.A. duales l'une de l'autre....)

\figure{\begin{tikzpicture}
\draw (0,0)--(1.80em,-0.025em);
\draw (.63em,-.29em)--(1.63em,-.29em);
\end{tikzpicture}}

\deleted{\unsure{C)}} Vers la fin des années 50\inserted{(ou début des 60 ?)}, alors que les grands fourbis cohomologiques ($\unsure{f_{!}|f^{!}},f^{!}$, coh. étale etc) sortaient à peine des limbes, le cours de Serre sur la théorie de Rosenlicht et Lang des Jacobiennes
% archive sheet number: 154
\scanpage{155}\mspage{10\makebox[0pt][r]{\tikz[baseline=0pt,overlay]{\draw (0.16em,0.70em) .. controls (0.16em,-0.37em) and (-1.16em,-0.39em) .. (-1.16em,0.44em);}}}
% The lower arc round 10 survives; its top meets the border of the copy.
généralisées, \deleted{et plu\unsure{s}}\inserted{et le corps de classe ``\unsure{géomrique|géomérique|géométrique}''} (voir le livre de Serre), et plus tard sa théorie ``géométrique'' du corps de classe \uline{local} à coups de groupes \unsure{proalgébriques|pro.algébriques} (voir son article sur le sujet), m'avait fait réfléchir à des formulations cohomologiques de ces résultats\inserted{ et d'autres}, qui devaient être notamment de nature \inserted{``géom'' telle} que les résultats ``arithmétiques'' sur un corps de base (ou résiduel) $k$ \inserted{général (fini, p.ex)} se ramènent immédiatement par \unsure{descente|descent} au cas ``géométrique'' du corps de base $\bar{k}$. J'ai eu un échange de lettres avec Serre -- je ne sais si j'en retrouverai des copies, mais je me rappelle qu'il s'agissait de projets de théories assez ambitieux, \inserted{à coups de résidus généralisés, Jacobiennes gén. locales etc,} dans au moins trois directions différentes. Mais je ne suis jamais parvenu, malgré de nombreux essais, à accrocher une bonne volonté à un de ces programmes. Voici quelques mots sur eux :

C) Situation $X$ de type fini sur \uline{corps} $k$,

\noindent\begin{tabular}{@{}|p{0.94\linewidth}@{}}
Construction d'un complexe de jacobiennes généralisées $J_{*X/k}$ (de long. égale à $\dim X$)
\end{tabular}
% archive sheet number: 155
\scanpage{156}\mspage{11}
C'était un complexe de pro-groupes algébriques sur $k$, affines à l'exception de $J_0$ si je me rappelle bien, (\deleted{\unsure{qui}}$J_0$ avait comme partie abélienne la partie abélienne de $\underline{\mathrm{Alb}}_{X/k}$
% margin note written sideways
\margin{\begin{tabular}{@{}|p{\dimexpr\marginparwidth-0.8em\relax}@{}}
\illegible{7} de $X$, i.e. on avait un \deleted{\unsure{f}}\unsure{hom} $X\longrightarrow\underline{H}_0$ universel pour les hom de $X$ dans des groupes loc. proalgébriques. Pour $X$ connexe, $\underline{H}_0$ est une ext. de $\mathbb{Z}$ par un groupe proalg. connexe
\end{tabular}}
% pen mark joining the note to the identification
\tikz[baseline=0pt]{\draw (0,.95ex)--(-.25em,.45ex)--(-.06em,.55ex)--(-.22em,.10ex)--(-.07em,.19ex)--(-.15em,-.14ex);}la jacobienne généralisée ordinaire). Sa construction, inspirée du complexe résiduel, passait par des jacobiennes généralisées (en un sens \inserted{cohomologique} convenable) des localisés $\operatorname{Spec}\underline{O}_{X,x}$ de $X$ en ses différents points. Il est possible qu'en un premier temps, je me limitais au cas $X$ lisse. Le rôle cohomologique de ce complexe était celui d'un complexe d'\uline{homologie}
\[
\mathrm{H}^{i}(X,G_X)\simeq\mathbb{E}\mathrm{xt}^{i}(J_{*X/k},G)\tag*{\makebox[0pt][r]{\makebox[\linewidth][l]{\unsure{\ensuremath{\alpha})|a)|(*)}}}}
\]
mais pour quels coefficients ? Je pense que je devais \deleted{\unsure{me limiter,|me limiter} plus ou moins à des} \inserted{prendre des groupes alg. comm. $G$ quelconques mais travailler en top. de Zariski (malédiction !)\unsure{, --|. --}} \deleted{$G$ connexes (brr !), ou \unsure{plu tôt|plutôt}} \inserted{même} dans le cas $G$ discret, je prenais néanmoins des $\mathrm{H}^{i}$ zariskiens -- cela donne des groupes de cohomologie un peu bêta, évidemment. Je réalisais qu'il faudrait travailler en coh etale ultérieurement, et que la construction\inserted{\unsure{$\times$|$x$}} des $J_{\deleted{\unsure{*}}iX/k}$ en serait évidemment modifiée. Quant à la sig%
% archive sheet number: 156
\scanpage{157}\mspage{12}%
nification des $\mathbb{E}\mathrm{xt}^{i}$ (hyperext), à un moment où Serre avait développé le formalisme des groupes proalgébriques, on n'avait pas trop peur de le prendre là-dedans -- au sens d'une catégorie dérivée qui à ce moment là n'avait encore jamais été définie et étudiée en forme.... (On a fait quand même un peu de chemin depuis !) J'ai bien l'impression, à la suite de ces antiques cogitations, pour heuristiques qu'elles \unsure{furent|fussent}, qu'on doit pouvoir développer à l'heure actuelle une telle théorie des $J_{*X/k}$, en coh fppf disons, donnant une formule $(*)$ sans limitation sur la dim $i$ de la cohomologie. (NB mais $J_*$ ne s'arrête évidemment plus en $\dim X=n$, mais en dim $2n$. Il est possible néanmoins que les composants $J_i$ soient de dim $0$ pour $\unsure{0|i|v}>n$).

Je crois bien que la construction de $J_*$ ne commute pas au chgt de base, mais que dans la cat. dérivée seulement elle y commute.

D) Soit $\deleted{f:}X/k$ schéma lisse \deleted{\unsure{si}} (pour simplifier) sur corps $k$, séparé de \unsure{T.f.|V.f.}, \inserted{de dim. relative $d$,} et $n$ entier $>0$. Si $n$ est premier à la caract. et si $F$ est un
% archive sheet number: 157
\scanpage{158}\mspage{13}
faisceau de coeff. sur $X$ annulé par $n$, la dualité globale nous dit que $Rf_{\deleted{\unsure{*}}!}(F)$ et $Rf_*(R\underline{\mathrm{Hom}}\,(F,\mu_n^{\otimes d}))$ sont duals l'un de l'autre à valeurs dans $(\mathbb{Z}/n\mathbb{Z})_k$, p.ex. $Rf_!(\unsure{Z|\mathbb{Z}}/n\mathbb{Z})$ et $Rf_*(\mu_n^{\otimes d})$, ou $Rf_!(\mu_n^{\otimes d})$ et $Rf_*(\mathbb{Z}/n\mathbb{Z})$, sont duals l'un de l'autre -- du moins modulo un shift d'amplitude $2d$ sur les dimensions. (Comme $\mathbb{Z}/n\mathbb{Z}$ est injectif sur lui-même, cela donne d'ailleurs des dualités parfaites
\[
R^{i}f_{\deleted{*}!}(F)\times R^{2d-i}f_*(\mathbb{R}\underline{\mathrm{Hom}}\,(F,\mu_n^{\otimes d}))\longrightarrow\mathbb{Z}/n\mathbb{Z}\big{)}.
\]

Si maintenant $n$ n'est plus supposé premier à la caractéristique, \inserted{p.ex \inserted{est} une puissance de $p=\mathrm{car}\,k>0$,} il semble que tout s'écroule, à commencer qu'on ne sait plus \inserted{(pour $d>1$)} par quoi remplacer $\mu_n^{\otimes d}$... Le miracle assez extraordinaire, c'est que pour $d=1$ i.e. $X$ une courbe lisse, tout continue à marcher parfaitement, pourvu qu'on l'énonce avec soin ! Les premières vérifications se font p.ex. avec $F=\mathbb{Z}/p\mathbb{Z}$, $\mu_p$, $\alpha_p$, $X$ complète -- on trouve que c'est OK grâce essentiellement à l'autodualité de la jacobienne. On peut sophistiquer les exemples, en prenant des coefficients $F$ \uline{tordus}, et $X$ non complète -- on se convainc que ça marche toujours.! Simplement il faut faire attention qu'ici les ${\unsure{R|\mathbb{R}}}^if_*(F)$, $R^if_!(F')$ ont une structure ``continue'' (ce sont essentiellement des groupes ind-pro-algébriques)
% archive sheet number: 158
\scanpage{159}\mspage{14}
Cela correspond au phénomène bien connu dans la théorie du corps de classe que la structure du $\pi_{1\,\mathrm{ab}}$ de $X$, quand $X$ est non complète, est \uline{continue} -- donc il en est de \unsure{$\bar{\mathrm m}$|$\bar{\mathrm n}$} du $\mathrm{H}^{1}(X,\mathbb{Z}/p^{n}\mathbb{Z})$ disons.

À ce propos, je te signale que Serre avait proposé (sans jamais l'écrire, je crois) une théorie de dualité des groupes algébriques commutatifs unipotents\inserted{ connexes}, \deleted{\unsure{duals|doubles}} pris \uline{mod isogénie radicielle}, dualité à valeurs dans $\mathbb{Q}/\mathbb{Z}$ (ou $\mathbb{Q}_p/\mathbb{Z}_p$). Il trouve que si (sur $k$ alg clos, disons) $G$ est un tel groupe, alors $G'=\operatorname{Ext}^{1}(G,\mathbb{Q}/\mathbb{Z})$ peut être muni \inserted{canoniquement} d'une structure de groupe quasi-algébrique (i.e. défini mod isog. radicielle), de façon unique sans doute en vérifiant qques propriétés de fonctorialité, et en exigeant que pour $G=\mathbb{G}_a$ on trouve sur $\operatorname{Ext}^{1}(\unsure{G_a|\mathbb{G}_a},\mathbb{Q}/\mathbb{Z})\simeq\mathbb{G}_a$ la structure habituelle. Soit $\Delta G=G'=\underline{\mathrm{Ext}}^{1}(G,\mathbb{Q}/\mathbb{Z})$, On trouve $G\simeq\Delta\Delta G$ i.e. $\Delta$ est une authentique autodualité ! J'appelle $\Delta$ la \uline{dualité de Serre}. Elle passe bien sûr aux ind-pro-groupes sur un corps de base quelconque (pas nec alg. clos) \inserted{de car $p>0$.} D'ailleurs, pour des groupes finis étale, c'est $\underline{\mathrm{Ext}}^{0}(G,\mathbb{Q}/\mathbb{Z})$ (dualité de \unsure{Pontryagin|Pontrjagin}) qui fournit une dualité parfaite. On peut \unsure{viser|visser} ensemble, dans une catégorie dérivée convenable, dualité de Serre et dualité de Pontryagin, en prenant
% archive sheet number: 159
\scanpage{160}\mspage{15}
$G\longmapsto\Delta G=\mathbb{R}\underline{\mathrm{Hom}}\,(G,\mathbb{Q}/\mathbb{Z})$ : on \unsure{l'appellera|l'appelle} encore dualité de Serre (``cohomologique''). Ce sera une magnifique autodualité, si on se place dans une catégorie dérivée où les $\underline{H}^{i}$ des complexes envisagés sont (à des passages \unsure{à|;} la limite près) des extensions de groupes étales par des groupes unipotents connexes. Or on ne rencontre que de tels complexes, en ``intégrant'' des coeff. finis $F$ sur $X$ par $Rf_!$ ou $Rf_*$. Ceci posé, passant à la limite sur $n$ dans la formulation initiale (ou encore partant avec $(\mathbb{Q}/\mathbb{Z})_k$ sur $k$, au lieu des $(\mathbb{Z}/n\mathbb{Z})_k$, et formant $f^!(\mathbb{Q}/\mathbb{Z})_k=(\mu_\infty)_X$), la formule de dualité prend la forme
\diagram{\[
\begin{tikzcd}[ampersand replacement=\&,column sep=0.8em,row sep=0.1em]
\& \text{``shift'' des dimensions} \\
\boxed{\Delta(Rf_!(F))\simeq Rf_*(DF[2])} \arrow[ur,start anchor=east,end anchor=west,shorten <=-0.8em] \&
\end{tikzcd}
\]}
\unsure{où} $D$ est la ``dualité de Cartier'' $R\underline{\mathrm{Hom}}\,(F,\mu_\infty)$ (ou $R\underline{\mathrm{Hom}}\,(F,\mathbb{G}_m)$ au choix ?), et $\Delta$ la dualité de Serre : cohomologie à support propre et à support quelconque s'échangent\inserted{ par dualité}, quand on prend ``en haut'' la dualité de Cartier, ``en bas'' la dualité de Serre.

La validité de la formule de dualité ne fait aucun doute -- le travail principal pour
% archive sheet number: 160
\scanpage{161}\mspage{16}
l'établir consiste certainement dans une description soigneuse des catégories de coefficients avec lesquels on travaille, aussi bien sur $X$ et sur $k$, et des foncteurs $\deleted{\unsure{\Delta}}D$ et $\deleted{\unsure{D}}\Delta$ dessus. Comme la définition d'une flèche est immédiate, une fois le travail de montage de la machine fait, la validité de la formule doit résulter sans difficulté des dévissages habituels, permettant de nous ramener au cas standart $F=\mathbb{Z}/p\mathbb{Z}$, $\mu_p$, $\alpha_p$ sur $X$ lisse complète (\uline{NB} le cas des coeff. premiers à la car. est déjà connu.)

% The opening of the following passage carries a stepped paragraph mark.
\noindent\tikz[baseline=-0.1ex,overlay]{\draw (-.43,0)--(-.17,0)--(-.17,.37)--(.53,.37);}Explicitant ce que donne la formule de dualité pour $R^1f_*(G_X)$, \unsure{où} $G$ est un groupe \inserted{fini} étale sur $k$ (le cas le plus important étant $G=(\mathbb{Z}/p^{\unsure{n|m}}\mathbb{Z})_k$), on retrouve la description à la Serre du \deleted{\unsure{``}}corps de classe ``géométrique'', en termes d'extensions par $G$ d'une jacobienne généralisée de $X$. \unsure{Ainsi}, la formule de dualité peut être comprise comme une version cohomologique considérablement enrichie du corps de \unsure{classes|classe} géométrique. Quand le corps de base $k$ est fini, pour retrouver le corps de classe sous sa forme classique, on peut mettre ``l'astuce de Lang'' \tikz[baseline=0pt]{\draw (.20em,0.78em)--(0,.78em);\draw (0,0.69em)--(0,-.14em)--(.20em,-.10em);}sur la \unsure{reltion|relation} entre \deleted{\unsure{un}}, le $\pi_1$ ``arithmétique'' d'un groupe algébrique \inserted{comm.} lisse connexe \inserted{$J$} sur $k$ et son $\mathrm{H}^{0}(k,J)=J(k)$, -- le $\pi_1^{\mathrm{ar}}(J)$ \unsure{classifiant} les isogénies \inserted{au dessus de $J$} à noyau un groupe \unsure{con}
% The last word reaches the bottom edge of the manuscript copy.
% archive sheet number: 161
